% !TEX program = xelatex
\documentclass[UTF8,12pt,a4paper]{ctexart}

\usepackage[margin=2.5cm]{geometry}
\usepackage{amsmath,amssymb,booktabs,array,multirow}
\usepackage{graphicx,xcolor,hyperref}
\usepackage{enumitem}
\usepackage{caption}
\captionsetup{font=small,labelfont=bf}
\hypersetup{colorlinks=true,linkcolor=blue,urlcolor=blue,citecolor=blue}
\setlength{\parindent}{2em}
\setlength{\parskip}{0.25em}

\newcommand{\todo}[1]{\textcolor{red!70!black}{【待补：#1】}}
\newcommand{\method}{Restormer 占位基线}

\title{从 CASSI 压缩测量直接预测多种植被指数：\\测量域任务建模与初步基线}
\author{作者信息待补}
\date{教师讨论稿（2026年6月29日）}

\begin{document}
\maketitle

\begin{abstract}
高光谱植被指数能够从不同波段组合中表征叶绿素、类胡萝卜素、生物量和环境胁迫等信息，
但传统处理流程通常需要先从编码孔径快照光谱成像（CASSI）测量中重构完整高光谱立方体，
再逐一计算指数。该两阶段流程将下游任务并不需要的全部光谱信息也纳入恢复目标，带来额外的
计算、存储和误差传播。本文研究从二维 CASSI 压缩测量直接预测 32 种植被指数的测量域任务。
我们建立统一的数据、归一化、训练、整叶测试和可视化流程，并以 Restormer 结构作为当前占位
基线验证任务可行性。现有基线已完成 300 轮训练，但旧测试程序存在重复归一化问题，故本文
暂不报告定量测试结论。当前稿件的目标是固定问题定义、论文逻辑和实验接口，为后续
Mamba、双分支结构及公式引导模型的统一比较提供基础。
\end{abstract}

\noindent\textbf{关键词：} CASSI；高光谱成像；植被指数；测量域计算；压缩任务；深度学习

\section{引言}

植被指数通过组合特定波段的反射率，能够突出与植被生理状态相关的光谱响应，在精准农业、
作物长势监测、叶绿素估计、胁迫识别和生态调查中具有重要作用。与只依赖少量宽波段的传统
多光谱观测相比，高光谱数据提供连续、细粒度的光谱信息，使同一场景可以计算多种具有不同
物理含义的指数。然而，高光谱数据的采集、传输与处理成本较高，限制了它在动态场景和资源
受限平台上的实时应用。\todo{补充植被指数与精准农业的权威综述引文}

CASSI 利用编码孔径和色散器件，将三维空间--光谱信息压缩到单帧二维测量中，因而能够避免
逐点或逐线扫描带来的时间延迟。现有 CASSI 研究主要围绕高光谱重构展开：先从测量恢复完整
高光谱立方体，再根据恢复波段计算 NDVI、SAVI 等指数。这一“重构后计算”范式把完整光谱
重建作为中间目标，但最终指数通常只依赖有限波段组合；重构误差还会经过差值、比值和倒数
运算传播到下游指数。由此产生的问题是：为了完成指数预测，是否必须恢复全部高光谱信息？
\todo{补充 CASSI 原理和代表性重构方法引文}

从测量域直接执行下游任务为上述问题提供了另一条路径。压缩任务方法尝试从编码测量中提取
与识别、检测或参数估计直接相关的判别信息，从而避免把资源消耗在完整信号恢复上。该思路
已经在视频快照压缩成像等场景中得到探索，但面向 CASSI 测量的稠密、多输出植被指数预测
仍缺少统一的问题定义与评估流程。特别是，不同植被指数的数值范围差异可跨越多个数量级，
若训练与测试归一化不一致，容易得到看似完整但无法比较的结果。\todo{补充测量域任务文献}

基于此，本文将 CASSI 测量直接映射为 32 通道植被指数图。当前工作首先固定数据接口、成像
模型、通道归一化和整叶评估协议，再以 Restormer 作为占位基线，为后续候选结构提供同一
比较起点。本文当前阶段的贡献可概括为：
\begin{enumerate}[leftmargin=2.5em]
  \item 建立从单帧 CASSI 压缩测量直接预测多种植被指数的统一任务定义，明确其与传统
  “重构--计算”流程的差异；
  \item 整理 84 波段输入、32 指数标签、固定物理掩模和整叶拼接评估的数据流程，并统一
  训练、测试与反归一化口径；
  \item 给出可复现的占位基线和候选模型接口，同时明确区分训练证据、有效测试指标与
  尚待验证的模型假设，避免使用失真的旧结果支撑结论。
\end{enumerate}

\section{相关工作}

\subsection{高光谱重构}
传统重构方法通常利用稀疏性、低秩性、总变分或非局部相似性构建优化问题，具有较强的物理
解释性，但迭代求解成本较高。深度学习方法通过端到端网络、即插即用先验或深度展开结构提高
恢复速度与质量，其中 Transformer 能够建模远距离空间--光谱依赖，Mamba 类状态空间模型则
尝试以更低的序列复杂度获得全局感受野。上述方法以恢复高保真光谱立方体为目标，而本文更
关心测量中与植被指数直接相关的信息。\todo{补充 MST、Restormer、Mamba 等原始论文}

\subsection{测量域任务}
测量域任务不要求显式恢复完整信号，而是联合利用成像编码和任务损失学习紧凑表示。与分类
或检测不同，植被指数预测需要为每个空间位置同时输出 32 个连续值，因此既要保持边缘和纹理，
又要处理指数间差异显著的量纲与动态范围。本文采用固定的逐通道上下界把标签映射到统一空间，
定量指标也在该空间计算。

\subsection{植被指数估计}
NDVI、SAVI、EVI、红边指数及色素相关指数分别利用红光、近红外、蓝光或红边波段的差值、
比值与倒数关系。本文使用 32 种指数，覆盖植被活力、叶绿素、类胡萝卜素、生物量及胁迫
响应。各指数公式与使用波段将在 `algorithm.md` 和附录中持续核验，当前稿不对尚未核验的
应用结论作强断言。

\section{问题定义}

设高光谱图像为 $\mathbf{X}\in\mathbb{R}^{H\times W\times L}$，其中 $L=84$，
波长范围为 445--860\,nm，采样间隔为 5\,nm。二维编码孔径为
$\mathbf{M}\in\mathbb{R}^{H\times W}$，色散步长为 $d=2$。简化的 CASSI 前向过程为
\begin{equation}
  \mathbf{Y}(i,j)=\frac{s}{L}\sum_{\ell=1}^{L}
  \mathbf{M}(i,j-d(\ell-1))\,
  \mathbf{X}(i,j-d(\ell-1),\ell),
  \label{eq:cassi}
\end{equation}
其中 $s=0.9$ 为当前实现中的测量缩放系数，越界位置取零。对 $256\times256$ patch，
测量宽度为 $256+(84-1)\times2=422$。

第 $k$ 个真实植被指数由相应波段函数 $g_k$ 产生：
\begin{equation}
  \mathbf{I}_k=g_k(\mathbf{X}),\qquad k=1,\ldots,32.
\end{equation}
传统流程先估计 $\hat{\mathbf{X}}$，再计算 $g_k(\hat{\mathbf{X}})$；本文直接学习
\begin{equation}
  \hat{\mathbf{I}}=f_{\theta}(\mathbf{Y},\mathbf{M}),
  \qquad \hat{\mathbf{I}}\in\mathbb{R}^{32\times H\times W}.
  \label{eq:direct}
\end{equation}

\section{占位方法与统一流程}

\subsection{测量展开与特征编码}
当前基线将二维测量按各波段对应的色散偏移截取为 84 个对齐通道，形成初始特征
$\mathbf{X}_0$。该操作不是高光谱重构，而是为网络提供与 CASSI 色散结构一致的输入表示。
随后通过卷积嵌入到基础特征维度。

\subsection{Restormer 占位基线}
\method 使用多尺度编码器--解码器。每个尺度包含通道注意力和门控深度卷积前馈模块，并将
84 通道掩模投影到当前特征维度参与注意力调制。解码器通过跳跃连接恢复空间细节，末端卷积
直接输出 32 个归一化指数通道。该结构仅用于建立当前可运行参照，不作为最终创新声明。

\subsection{训练目标}
标签按每个指数预设的最小值和最大值归一化。当前训练目标由 Charbonnier、$L_1$ 与结构
相似性损失组成：
\begin{equation}
  \mathcal{L}=\mathcal{L}_{\mathrm{char}}+
  \mathcal{L}_{1}+\mathcal{L}_{\mathrm{SSIM}}.
\end{equation}
测试时模型输出直接与归一化标签比较，不再对输出重复执行固定 min--max。

\begin{figure}[htbp]
  \centering
  \fbox{\parbox[c][4.2cm][c]{0.88\textwidth}{\centering
  图示占位：HSI 与 mask 经过 CASSI 前向模型形成二维测量；传统流程先重构再计算指数，
  本文流程由网络直接输出 32 通道指数。}}
  \caption{传统重构--计算流程与本文测量域直接预测流程。}
  \label{fig:pipeline}
\end{figure}

\section{实验设置与初步状态}

\subsection{数据与划分}
转换后的数据包含连续编号的 252 组 HSI 和指数标签，以及一个固定 mask。计划中的活动划分为
1--200 号训练、201--252 号测试。现有 Restormer 权重训练时仍沿用了原始 MAT 阶段的
`skip\_ids=[15,43,109]`，故实际使用 197 个训练源场景，每轮随机采样 2560 个
$256\times256$ patch。本文保留该旧配置快照，后续模型将使用连续编号划分并统一重评。

\subsection{训练配置}
占位基线基础维度为 48，四层块数为 $[4,6,6,8]$，注意力头数为 $[1,2,4,8]$。
训练采用 AdamW，初始学习率 $10^{-4}$，batch size 为 1，共 300 epoch；梯度裁剪阈值为
0.2。输入测量逐样本归一化，标签使用固定逐通道归一化。

\subsection{当前可用证据}
\begin{table}[htbp]
  \centering
  \caption{Restormer 占位基线的训练日志（非测试集性能）。}
  \begin{tabular}{ccc}
    \toprule
    Epoch & 训练 loss & 训练 PSNR / dB \\
    \midrule
    50  & 0.13149128 & 35.14 \\
    100 & 0.12659922 & 35.86 \\
    150 & 0.12311621 & 36.24 \\
    200 & 0.12143392 & 36.45 \\
    250 & 0.11937794 & 36.76 \\
    300 & 0.11846657 & 36.85 \\
    \bottomrule
  \end{tabular}
\end{table}

旧测试程序存在重复归一化，相关数值已判定为不可用于论文。修订后统一测试表如下保留位置：
\begin{table}[htbp]
  \centering
  \caption{统一测试协议下的结果占位。}
  \begin{tabular}{lcccc}
    \toprule
    模型 & 平均 $L_1$ & 平均 PSNR & SSIM & 状态 \\
    \midrule
    Restormer & -- & -- & -- & 待 GPU 空闲后重评 \\
    MST-Mamba & -- & -- & -- & 待训练 \\
    WPO3D & -- & -- & -- & 待训练 \\
    DHM & -- & -- & -- & 待训练 \\
    IFGNet & -- & -- & -- & 待训练 \\
    \bottomrule
  \end{tabular}
\end{table}

\section{讨论与后续模型}
当前结果只能说明训练流程能够收敛，不能说明模型在独立测试场景上的精度，也不能据此比较
候选结构。后续实验将重点回答三类问题：Mamba 的线性序列建模能否降低大尺寸特征的开销；
空间--光谱或全局--局部双分支能否缓解单一扫描的局部信息损失；显式重建少量关键波段并
利用指数公式约束，能否提高物理一致性。所有方案必须共享数据划分、归一化、checkpoint
选择和整叶评估协议。

对于 CRI、ARI 等含倒数的指数，还应单独分析低反射率分母导致的长尾和数值稳定性；对于
不同量纲指数，仅报告原量纲平均误差会被大范围通道支配，因此主要比较应在统一归一化空间
进行，并同时给出逐指数结果。

\section{结论}
本文建立了从 CASSI 二维压缩测量直接预测 32 种植被指数的任务框架，整理了成像、归一化、
训练、整叶测试和可视化接口，并以 Restormer 作为占位基线。当前训练日志提供了流程收敛的
初步证据，但由于旧测试重复归一化，本文主动暂缓定量测试结论。下一阶段将在修订后的统一
协议下重新评估占位权重，并完成 Mamba、双分支和公式引导模型的公平比较。

\appendix
\section{当前待补材料}
\begin{itemize}
  \item 权威引用与完整 BibTeX；
  \item 修订后测试的 32 指数逐通道表格；
  \item 完整叶片预测、标签和误差图；
  \item 模型参数量、FLOPs、推理时间和显存比较；
  \item 最终方法消融实验与失败案例分析。
\end{itemize}

\end{document}
